\documentclass[french]{article}
\usepackage{verbatim}
\usepackage{amsmath,amsfonts,amssymb}
\usepackage{pgfplots}
\usepackage[utf8]{inputenc}
\usepackage[french]{babel}
\usepackage{pgf,tikz,tkz-tab}
\usepackage{mathrsfs}
\usetikzlibrary{arrows}
\usepackage{pstricks}
\usepackage{algorithm}
\usepackage{algorithmic}
%\pgfplotsset{compat=1.15}
\usepackage{geometry}
\geometry{hmargin=0.5cm,vmargin=0.7cm}


\begin{document}
\twocolumn
\begin{center}
\underline{\textbf{Devoir à la maison II}}
\end{center}

\section*{Exercice I}
\begin{itemize}
\item[A)] Calculer $ 1 \times 2  \times 3 $.
\item[B)] Calculer $ 1 \times 2  \times 3  \times 4$.\\

\end{itemize}

On appelle $3!$ ( se prononce 3 factoriel ) le nombre $ 1 \times 2  \times 3 $ puis $4!$ le nombre $ 1 \times 2  \times 3  \times 4$ etc.

\begin{itemize}
\item[C)] Completer l'égalité suivante $ 4! = 3! \times ... $\\

\end{itemize}

\underline{Recopier et completer les phrases suivante :}

$4!$ est un ... de $3!$.

$3!$ est un ... de $4!$.\\


\begin{itemize}
\item[D)] Calculer $6!$.
\item[E)] Réduire $\frac{5!}{4!}$.
\item[F)] Décomposer $4!$ en produit de nombre premier.
\item[G)] 5 est-il un nombre premier?
\item[H)] Décomposer $5!$ en produit de nombre premier \underline{sans calcul} à l'aide de la question E et F.\\
\end{itemize}

Il est difficile de décomposer en produit de nombre premier un grand nombre, surtout quand celui çi est le résultat d'une multiplication par des grands nombres premiers.\\

Essayer de décomposer en nombre premier à la main, le nombre $10403$.  Puis aller sur https://www.isprimenumber.com, et taper le nombre $10403$,le site vous donnera la décomposition en nombre premier de celui-çi.

\begin{center}
\underline{\textbf{Devoir à la maison II}}
\end{center}

\section*{Exercice I}
\begin{itemize}
\item[A)] Calculer $ 1 \times 2  \times 3 $.
\item[B)] Calculer $ 1 \times 2  \times 3  \times 4$.\\

\end{itemize}

On appelle $3!$ ( se prononce 3 factoriel ) le nombre $ 1 \times 2  \times 3 $ puis $4!$ le nombre $ 1 \times 2  \times 3  \times 4$ etc.

\begin{itemize}
\item[C)] Completer l'égalité suivante $ 4! = 3! \times ... $\\

\end{itemize}

\underline{Recopier et completer les phrases suivante :}

$4!$ est un ... de $3!$.

$3!$ est un ... de $4!$.\\


\begin{itemize}
\item[D)] Calculer $6!$.
\item[E)] Réduire $\frac{5!}{4!}$.
\item[F)] Décomposer $4!$ en produit de nombre premier.
\item[G)] 5 est-il un nombre premier?
\item[H)] Décomposer $5!$ en produit de nombre premier \underline{sans calcul} à l'aide de la question E et F.\\
\end{itemize}

Il est difficile de décomposer en produit de nombre premier un grand nombre, surtout quand celui çi est le résultat d'une multiplication par des grands nombres premiers.\\

Essayer de décomposer en nombre premier à la main, le nombre $10403$.  Puis aller sur https://www.isprimenumber.com, et taper le nombre $10403$,le site vous donnera la décomposition en nombre premier de celui-çi.\\


\begin{center}
\underline{\textbf{Devoir à la maison II}}
\end{center}
\section*{Exercice I}
\begin{itemize}
\item[A)] Calculer $ 1 \times 2  \times 3 $.
\item[B)] Calculer $ 1 \times 2  \times 3  \times 4$.\\

\end{itemize}

On appelle $3!$ ( se prononce 3 factoriel ) le nombre $ 1 \times 2  \times 3 $ puis $4!$ le nombre $ 1 \times 2  \times 3  \times 4$ etc.

\begin{itemize}
\item[C)] Completer l'égalité suivante $ 4! = 3! \times ... $\\

\end{itemize}

\underline{Recopier et completer les phrases suivante :}

$4!$ est un ... de $3!$.

$3!$ est un ... de $4!$.\\


\begin{itemize}
\item[D)] Calculer $6!$.
\item[E)] Réduire $\frac{5!}{4!}$.
\item[F)] Décomposer $4!$ en produit de nombre premier.
\item[G)] 5 est-il un nombre premier?
\item[H)] Décomposer $5!$ en produit de nombre premier \underline{sans calcul} à l'aide de la question E et F.\\
\end{itemize}

Il est difficile de décomposer en produit de nombre premier un grand nombre, surtout quand celui çi est le résultat d'une multiplication par des grands nombres premiers.\\

Essayer de décomposer en nombre premier à la main, le nombre $10403$.  Puis aller sur https://www.isprimenumber.com, et taper le nombre $10403$,le site vous donnera la décomposition en nombre premier de celui-çi.\\

\begin{center}
\underline{\textbf{Devoir à la maison II}}
\end{center}


\section*{Exercice I}
\begin{itemize}
\item[A)] Calculer $ 1 \times 2  \times 3 $.
\item[B)] Calculer $ 1 \times 2  \times 3  \times 4$.\\

\end{itemize}

On appelle $3!$ ( se prononce 3 factoriel ) le nombre $ 1 \times 2  \times 3 $ puis $4!$ le nombre $ 1 \times 2  \times 3  \times 4$ etc.

\begin{itemize}
\item[C)] Completer l'égalité suivante $ 4! = 3! \times ... $\\

\end{itemize}

\underline{Recopier et completer les phrases suivante :}

$4!$ est un ... de $3!$.

$3!$ est un ... de $4!$.\\


\begin{itemize}
\item[D)] Calculer $6!$.
\item[E)] Réduire $\frac{5!}{4!}$.
\item[F)] Décomposer $4!$ en produit de nombre premier.
\item[G)] 5 est-il un nombre premier?
\item[H)] Décomposer $5!$ en produit de nombre premier \underline{sans calcul} à l'aide de la question E et F.\\
\end{itemize}

Il est difficile de décomposer en produit de nombre premier un grand nombre, surtout quand celui çi est le résultat d'une multiplication par des grands nombres premiers.\\

Essayer de décomposer en nombre premier à la main, le nombre $10403$.  Puis aller sur https://www.isprimenumber.com, et taper le nombre $10403$,le site vous donnera la décomposition en nombre premier de celui-çi.
\section*{Exercice II}

On sait qu'il est difficile de savoir quand un nombre est premier ou non, pour autant il n'est pas non plus impossible de le savoir.

Le but ici est donc de faire un algorithme pour savoir si un nombre est premier ou non.\\


\begin{itemize}
\item[A)] Rappeler ce qu'est un nombre premier.
\item[B)] Trouver une méthode pour déterminer si un nombre est premier ou non. \\

\end{itemize}

\underline{Recopier et compléter l'algorithme suivant :}

\begin{algorithm}
\caption{Test de primalité}
\begin{algorithmic} 
\REQUIRE $T>2$
\STATE $n \leftarrow 2$
\WHILE{$n < T$}
\STATE $R \leftarrow ...$
\IF{$R = 0$}
\STATE Afficher 'T ... premier'
\STATE $n \leftarrow T+1$
\ELSE
\STATE $n \leftarrow n+1$
\ENDIF
\ENDWHILE
\IF{n=T}
\STATE Afficher 'T ... premier'
\ENDIF
\end{algorithmic}
\end{algorithm}		

\vspace{4cm}
\end{comment}

\end{document}